\section{Rational orbital estimates and the ambient multiplicity}
\label{amb-section}

The dimension of the decorated, definite source space in
Section~\ref{cd-section} is different from the multiplicity in the
indefinite automorphic space. Here we isolate that second count,
prove its rational-element and distant-orbit estimates, and give the
exact analytic reduction needed for the final comparison. The
coefficients throughout this section are complex. In particular, the
large auxiliary prime below imposes a condition on the finite level;
it is not a coefficient characteristic. These are the conventions of
the submitted multiplicity theorem.\cite{alpo-uniform-source}

\subsection{The multiplicity and its proposed main term}

Let $E/F$ be a CM extension, with $F$ totally real of degree $f$,
$[E:F]=2$, and $E/F$ unramified at all finite places. Let $V$ be an
$n$-dimensional Hermitian space over $E$, with a self-dual
$\mathcal O_E$-lattice $\Lambda_0$. Write $G=\operatorname U(V)$,
$\Gamma=G(F)$ and
\[
 G_\infty=\prod_{v\mid\infty}G(F_v),\qquad
 G_f=G(\mathbb A_{F,f}),\qquad
 \mathcal X=\Gamma\backslash(G_\infty\times G_f).
\]
At every real place the signature is $(p,q)$, where $p+q=n$.
Choose coordinates in which the Hermitian matrix is a diagonal
matrix $J_v$ with entries $1$ and $-1$. No ordering of these signs
is assumed. The associated positive norm is the usual Euclidean
norm in these coordinates. Put
\[
                       r_E=|D_E|^{1/(2f)}.
\]
Let $K_f$ stabilize $\widehat\Lambda_0$, and let $K_f(g)$ be its
principal congruence subgroup of positive integer level $g$.

The finite test function $e:G_f\to\mathbb C$ is locally constant,
compactly supported in $K_f$, bi-invariant under $K_f(g)$, and
satisfies
\begin{equation}\label{amb-idempotent}
 e*e=e,\qquad e(y^{-1})=\overline{e(y)},\qquad
 \sum_{i,j}\overline{c_i}c_j e(y_i^{-1}y_j)\ge0.
\end{equation}
The last inequality holds for every finite collection of $y_i,c_i$.
There is a rational prime $P>4n$ which splits completely in $E$,
such that
\begin{equation}\label{amb-unipotent-level}
 e(y)\ne0\quad\Longrightarrow\quad
             (y-1)^n\widehat\Lambda_0\subseteq
                                      P\widehat\Lambda_0.
\end{equation}
Fix Haar measures $dg_v$ on $G(F_v)$ and $dg_f$ on $G_f$.

For each real place let $\sigma_v$ be the discrete-series
representation with the specified Harish--Chandra parameter
$\lambda_v$, and choose a unit vector $w_v$ in its lowest
$K_v$-type. Let $d_v>0$ be its formal degree for $dg_v$ and
$\sigma=\widehat\bigotimes_v\sigma_v$. With $R$ the right regular
representation on $L^2(\mathcal X)$, define
\begin{equation}\label{amb-D-definition}
 D=\dim_{\mathbb C}
   \bigl\{R(e)T:T\in
       \operatorname{Hom}^{\mathrm{cont}}_{G_\infty}
                         (\sigma,L^2(\mathcal X))\bigr\}.
\end{equation}
The dimension in \eqref{amb-D-definition} is used as a finite number
only after finite-dimensionality has been established. Write
\begin{equation}\label{amb-main-term}
 m_0=\Bigl(\prod_v d_v\Bigr)e(1),\qquad
 D_*=m_0\operatorname{vol}(\mathcal X).
\end{equation}
Lemma~\ref{amb-positive-type} below shows that $e(1)$ is real and
nonnegative. The source's identification of its row multiplicity
with \eqref{amb-D-definition} is the identification used to write
$D=D_{\rm all}$; it does not change $D_*$.

The claimed comparison is
\begin{equation}\label{amb-target}
                 |D-D_*|\le
                  \exp\bigl(-nf\log n/16\bigr)D_*.
\end{equation}
The source proves this for sufficiently large $n$, uniformly under
the following quantitative assumptions. Here $C_g\ge0$, $b\ge2$,
and $h$ is a lower bound for every difference
$|\lambda_{v,i}-\lambda_{v,j}|$ whose two indices have opposite
signs in $J_v$:
\begin{align}
 q^2&\le n,& \log f&\le\sqrt n,& \log r_E&\le e^n,
                                                        \notag\\
 |\lambda_{v,i}|&\le b^2,&1\le g&\le b^{C_g},&
 \log b&\le\exp(nf\log n/16),                         \label{amb-parameters}\\
 h&\ge8\cdot10^{14}(C_g+4)n^{19}f^5\log(br_E).
                                                        \notag
\end{align}
The size of the threshold for $n$ may depend on $C_g$.
The statements proved below specify exactly which additional
analytic inputs are still needed for \eqref{amb-target}.

\subsection{Unipotence of nearby rational isometries}

For an invertible complex matrix $A$ with singular values
$s_1,\ldots,s_n$, set
\[
             t(A)=\Bigl(\sum_i(\log s_i)^2\Bigr)^{1/2}.
\]
Its polar decomposition is $A=k\exp X$, where $k$ is unitary,
$X$ is Hermitian, and $\|X\|_{\rm HS}=t(A)$. In particular,
$\|A\|_{\rm op}\le e^{t(A)}$.

\begin{lemma}\label{amb-archimedean-traces}
Suppose $n\ge1$, $\epsilon\ge0$, $n\epsilon\le1/4$ and
$t(A)\le\epsilon$. Then, for $0\le j\le n$,
\begin{equation}\label{amb-power-traces}
                  |\operatorname{tr}(A^j)-n|\le3n.
\end{equation}
If $\operatorname{tr}A=n$, then
\begin{equation}\label{amb-close-to-one}
                  \|A-1\|_{\rm HS}\le3\sqrt{n\epsilon}.
\end{equation}
\end{lemma}
\begin{proof}
The operator-norm bound gives
$|\operatorname{tr}(A^j)|\le n\|A^j\|_{\rm op}
\le ne^{j\epsilon}\le ne^{1/4}<3n/2$.
The last strict bound follows, for example, by comparing the
exponential series with $\sum_{k\ge0}(1/4)^k=4/3$.
The triangle inequality proves \eqref{amb-power-traces}.
If the trace is $n$, then
\[
 \|A-1\|_{\rm HS}^2
   =\operatorname{tr}(A^*A)+n-2\Re\operatorname{tr}A
   =\operatorname{tr}(e^{2X})-n
   \le n(e^{2\epsilon}-1).
\]
Here $\epsilon\le1/4$. The exponential series gives
$e^{2\epsilon}-1\le2\epsilon/(1-2\epsilon)
\le4\epsilon$. Taking square roots proves the stronger bound
$2\sqrt{n\epsilon}$, hence \eqref{amb-close-to-one}.
\end{proof}

\begin{lemma}\label{amb-integral-traces}
Let $E$ be a number field, $\Lambda$ a full
$\mathcal O_E$-lattice in $E^n$, and $A\Lambda\subseteq\Lambda$.
Suppose that $P\mathcal O_E$ is radical and
$(A-1)^n\Lambda\subseteq P\Lambda$. Then
\[
               \operatorname{tr}(A^j)-n\in P\mathcal O_E
                                      \qquad(j\ge0).
\]
No freeness of $\Lambda$ over $\mathcal O_E$ is required.
\end{lemma}
\begin{proof}
At every nonzero prime $\mathfrak p$, the localization
$\Lambda_{\mathfrak p}$ is a free module of rank $n$ over the
discrete valuation ring $\mathcal O_{E,\mathfrak p}$.
Thus the characteristic polynomial of $A$ has coefficients in
every $\mathcal O_{E,\mathfrak p}$, hence in their intersection
$\mathcal O_E$. In particular all the traces in question are
integral. At a prime $\mathfrak p\mid P$, the reduction of $A$
on $\Lambda_{\mathfrak p}/\mathfrak p\Lambda_{\mathfrak p}$ is
unipotent. Its $j$th power has all eigenvalues equal to one over
an algebraic closure of the residue field, so its trace is $n$.
Consequently $\operatorname{tr}(A^j)-n$ belongs to every such
$\mathfrak p$. Radicality gives
$\bigcap_{\mathfrak p\mid P}\mathfrak p=P\mathcal O_E$.
\end{proof}

\begin{proposition}\label{amb-near-unipotent}
Let $\Lambda,A,P$ be as in Lemma~\ref{amb-integral-traces}, with
$P>3n$. For every complex embedding $\iota:E\hookrightarrow
\mathbb C$, suppose that $\iota(A)$ is similar to a matrix
$A_\iota$ satisfying $t(A_\iota)\le\epsilon$, where
$n\epsilon\le1/4$. Then
\begin{equation}\label{amb-unipotence}
                   \det(T-A)=(T-1)^n,
                   \qquad(A-1)^n=0.
\end{equation}
Moreover \eqref{amb-close-to-one} holds for each $A_\iota$.
\end{proposition}
\begin{proof}
For $1\le j\le n$ put $a_j=\operatorname{tr}(A^j)-n$.
Lemma~\ref{amb-integral-traces} writes $a_j=P b_j$ with
$b_j\in\mathcal O_E$. If $a_j\ne0$, then
\[
 |N_{E/\mathbb Q}(a_j)|
       =P^{[E:\mathbb Q]}|N_{E/\mathbb Q}(b_j)|
       \ge P^{[E:\mathbb Q]}.
\]
But similarity invariance of trace and
\eqref{amb-power-traces} bound the same norm by
$(3n)^{[E:\mathbb Q]}$, a contradiction. Thus
$\operatorname{tr}(A^j)=n$ for $1\le j\le n$.

For completeness these power traces determine the characteristic
polynomial in characteristic zero. If $\alpha_i$ are its roots
with multiplicity, then the polynomial
$Q(z)=\prod_i(1-\alpha_i z)$ satisfies the formal identity
\[
             -zQ'(z)=Q(z)\sum_{j\ge1}
                              \operatorname{tr}(A^j)z^j.
\]
Writing $Q(z)=\sum_{j=0}^n c_jz^j$ with $c_0=1$, comparison
of degree $j$ determines $j c_j$ from $c_0,\ldots,c_{j-1}$ and
the first $j$ traces. Those traces agree with the identity matrix,
so $Q(z)=(1-z)^n$. Cayley--Hamilton now gives
$(A-1)^n=0$. Finally the first trace is $n$, and
Lemma~\ref{amb-archimedean-traces} gives the asserted bound.
\end{proof}

For an adelic point $x=(x_\infty,x_f)$, define its rational
lattice and positive metrics by
\begin{equation}\label{amb-point-lattice}
 \Lambda_x=V\cap x_f\widehat\Lambda_0,
 \qquad\|u\|_{x,v}=\|x_v^{-1}u_v\|_2.
\end{equation}
Here $u_v$ is expressed in the chosen $J_v$-coordinates.
The lattice is full and self-dual. Indeed a common integer clears
the denominators of $x_f$ and $x_f^{-1}$, which bounds
$\Lambda_x$ between two homothetic copies of $\Lambda_0$.
The Chinese remainder theorem identifies its completions with
$x_f\widehat\Lambda_0$; duality can then be checked at each
finite prime. This also proves
\[
 (x^{-1}\gamma x)_f\in K_f\Longrightarrow
                        \gamma\Lambda_x=\Lambda_x,
\]
and \eqref{amb-unipotent-level} gives
$(\gamma-1)^n\Lambda_x\subseteq P\Lambda_x$ whenever the
finite test function is nonzero.

Complete splitting of $P$ makes $P\mathcal O_E$ radical. The
$f$ selected complex embeddings and their conjugates are all the
embeddings of $E$. Therefore Proposition~\ref{amb-near-unipotent}
applies to every rational summand with
$e((x^{-1}\gamma x)_f)\ne0$ and
$t((x^{-1}\gamma x)_v)\le\epsilon$ for all $v$.
In particular it applies at the exact radius
\begin{equation}\label{amb-small-radius}
                  \epsilon=10^{-6}n^{-8}f^{-2}.
\end{equation}
The source's further radius bounds follow directly: if
$\delta=3\sqrt{n\epsilon}$ and $q^2\le n$, then
\[
             \delta<\tfrac14,
             \qquad\delta qnf
               =3\cdot10^{-3}\frac q{n^{5/2}}
               \le3\cdot10^{-3}<\tfrac1{100}.
\]

\subsection{A separation bound from normalized covolumes}

We prove the separation estimate used in the orbit count without
requiring the construction of a Harder--Narasimhan filtration.
Put $c_E=2^{-f}|D_E|^{1/2}$. For a nonzero
$\mathcal O_E$-lattice $M$ of rank $r$, with its positive metrics
at the $f$ complex places, define
\begin{equation}\label{amb-size}
 H(M)=c_E^{-r}\operatorname{covol}_{\mathbb R}(M),\qquad
 s(M)=H(M)^{1/(2fr)},\qquad
 \lambda(\Lambda)=\inf_{0\ne M\subseteq\Lambda}s(M).
\end{equation}
Covolumes are measured in the real span of the lattice, with the
Euclidean metric which is the sum of the selected complex metrics.
For a torsion-free quotient, use the orthogonal quotient metrics.
These are exactly the normalized covolumes and exponents in the
submitted least-size function.

\begin{lemma}\label{amb-size-duality}
For every full $\mathcal O_E$-lattice $\Lambda$ with positive
metrics, $\lambda(\Lambda)>0$. Let $\Lambda^*$ be its
$\mathcal O_E$-linear dual, with the dual metrics. Then
\begin{equation}\label{amb-height-dual}
 H(\Lambda^*)=H(\Lambda)^{-1},\qquad
 \sup_{\substack{S\subsetneq\Lambda\\\Lambda/S\ \text{torsion-free}}}
              s(\Lambda/S)=\lambda(\Lambda^*)^{-1}.
\end{equation}
If $\Lambda$ is self-dual for a Hermitian form whose matrix is
$J_v$ in an orthonormal basis of each positive metric, then
$\lambda(\Lambda^*)=\lambda(\Lambda)$.
\end{lemma}
\begin{proof}
For $1\le r\le n$, a $\mathbb Z$-basis of any rank-$r$
$\mathcal O_E$-sublattice gives a nonzero vector in the fixed
lattice $\bigwedge^{2fr}_{\mathbb Z}\Lambda$. Its Euclidean
length is the covolume of the sublattice. A lattice has a positive
minimum length for nonzero vectors: an invertible real linear
map from $\mathbb Z^N$ bounds every such length below by its
smallest singular value. Taking the minimum over the finitely
many ranks proves positivity in \eqref{amb-size}.

If $M_0=\bigoplus_{i=1}^r\mathcal O_E b_i$ is free and
$G_v=(\langle b_i,b_j\rangle_v)$, the discriminant determinant
for the Minkowski embedding gives
\begin{equation}\label{amb-free-height}
                       H(M_0)=\prod_v\det G_v.
\end{equation}
Explicitly the covolume of $\mathcal O_E$ in its $f$ complex
coordinates is $c_E$; the real determinant of the complex change
of basis at $v$ is $|\det B_v|^2=\det G_v$. Multiplying these
determinants proves \eqref{amb-free-height}. The dual Gram
matrices are $G_v^{-1}$, so the dual-height assertion follows
for $M_0$.

For arbitrary $M$, choose such a free sublattice $M_0\subset M$
of finite index $a$. Then $H(M)=H(M_0)/a$ and
$[M_0^*:M^*]=a$. The latter index equality can be checked at
each prime: diagonalizing the inclusion over the discrete
valuation ring changes its exponents to their negatives in the
dual inclusion, with the same lengths of the two finite
quotients. Their global orders are the products of the residue
norms to these lengths. Thus
$H(M^*)=aH(M_0^*)=H(M)^{-1}$, without assuming $M$ is free.

If $S$ is saturated in $\Lambda$, then $Q=\Lambda/S$ is
projective over the Dedekind domain $\mathcal O_E$.
Dualizing the split module sequence identifies $Q^*$ with the
saturated submodule of $\Lambda^*$ annihilating $S$, with its
induced metric. Conversely every saturated submodule of
$\Lambda^*$ arises in this way, by dualizing its projective
quotient. Saturating a nonzero sublattice decreases its covolume
and preserves its rank; consequently the infimum in
\eqref{amb-size} is unchanged if it is restricted to saturated
sublattices. Since $s(Q^*)=s(Q)^{-1}$, this proves the supremum
identity.

Finally the Hermitian form identifies $\Lambda$ conjugate
linearly with $\Lambda^*$. At each complex place the norm of the
functional $u\mapsto\langle y,J_vu\rangle$ is $\|y\|$, because
$J_v$ is unitary. This identification is an isometry, takes
$\mathcal O_E$-sublattices to such sublattices and preserves their
real covolumes. Hence it preserves the infimum defining
$\lambda$.
\end{proof}

\begin{lemma}\label{amb-map-separation}
Under the self-duality hypothesis of Lemma~\ref{amb-size-duality},
every nonzero $E$-linear map $T$ with
$T\Lambda\subseteq\Lambda$ satisfies
\begin{equation}\label{amb-geometric-separation}
            \Bigl(\prod_v\|T_v\|_{\rm op}\Bigr)^{1/f}
                                     \ge\lambda(\Lambda)^2.
\end{equation}
\end{lemma}
\begin{proof}
Let $S=\Lambda\cap\ker T$, $Q=\Lambda/S$, and $M=T\Lambda$,
of common positive rank $r$. The first is saturated, and the
induced map $Q\to M$ is an isomorphism of
$\mathcal O_E$-modules. At each complex place its singular
values are at most $\|T_v\|_{\rm op}$: identify the quotient
with the orthogonal complement of $\ker T_v$ and restrict $T_v$.
The change-of-volume formula, applied first to a finite-index
free sublattice and then cancelling that common index, gives
\[
                 H(M)\le H(Q)\prod_v\|T_v\|_{\rm op}^{2r}.
\]
Take $2fr$th roots. By Lemma~\ref{amb-size-duality},
$s(M)\ge\lambda(\Lambda)$ and
$s(Q)\le\lambda(\Lambda)^{-1}$. These inequalities give
\eqref{amb-geometric-separation}.
\end{proof}

For the point lattice \eqref{amb-point-lattice}, put
\begin{equation}\label{amb-tau}
          \lambda_x=\lambda(\Lambda_x),\qquad
          \tau_x=\min(1/2,\lambda_x)>0.
\end{equation}
The positive metrics in \eqref{amb-point-lattice} are compatible
with the original Hermitian form, since $x_v$ is an isometry of
$J_v$. Thus Lemma~\ref{amb-map-separation} applies. If distinct
rational elements $\gamma,\eta$ both have
$(x^{-1}\gamma x)_f,(x^{-1}\eta x)_f\in K_f$, their difference
preserves $\Lambda_x$, and
\begin{equation}\label{amb-two-matrices}
 \Bigl(\prod_v\|(x^{-1}\gamma x)_v
                         -(x^{-1}\eta x)_v\|_{\rm op}\Bigr)^{1/f}
                                                    \ge\tau_x^2.
\end{equation}

\subsection{Counting rational matrices and summing the tail}

\begin{lemma}\label{amb-packing}
Let $W_1,\ldots,W_f$ be real normed spaces of dimension $N$,
where $f,N\ge1$. Suppose a set $S\subset\prod_vW_v$ satisfies
\[
       \Bigl(\prod_v\|s_v-t_v\|\Bigr)^{1/f}\ge a>0
                       \quad(s,t\in S,\ s\ne t).
\]
For positive $R_v$, the subset with $\|s_v\|\le R_v$ for all
$v$ is finite, of cardinality at most
\begin{equation}\label{amb-packing-bound}
              \left(1+\frac{2(\prod_vR_v)^{1/f}}a\right)^{Nf}.
\end{equation}
\end{lemma}
\begin{proof}
Rescale the $v$th coordinate by $R/R_v$, where
$R=(\prod_vR_v)^{1/f}$. The product of the scale factors is one,
so the separation hypothesis is unchanged. In the maximum norm
on $\prod_vW_v$, the new set is $a$-separated and lies in the
closed ball of radius $R$. The open balls of radius $a/2$ about
any finite selection of its points are disjoint and lie in the
ball of radius $R+a/2$. Lebesgue volume in dimension $Nf$ gives
\eqref{amb-packing-bound} for every finite selection. An infinite
set has finite subsets of arbitrarily large cardinality, so the
same bound proves finiteness as well.
\end{proof}

\begin{corollary}\label{amb-rational-ball}
For $r_v\ge0$, the rational elements with finite component in
$xK_fx^{-1}$ and $t((x^{-1}\gamma x)_v)\le r_v$ form a finite
set of cardinality at most
\begin{equation}\label{amb-ball-bound}
 A_x\exp\left(2n^2\sum_vr_v\right),\qquad
                   A_x=\left(\frac{3\sqrt n}{\tau_x^2}\right)^{2n^2f}.
\end{equation}
\end{corollary}
\begin{proof}
Use \eqref{amb-two-matrices} and Lemma~\ref{amb-packing} in the
spaces of complex $n\times n$ matrices with operator norm,
each of real dimension $2n^2$, with $a=\tau_x^2$ and
$R_v=e^{r_v}$. The map from rational elements to their matrices
is injective, since an embedding of $E$ is injective.
Writing $R=e^{\sum_vr_v/f}\ge1$, we have
$1+2R/\tau_x^2\le3R/\tau_x^2\le3\sqrt nR/\tau_x^2$.
Raise this inequality to $2n^2f$. The proof actually permits
omitting $\sqrt n$; we keep it to match the source bound.
\end{proof}

\begin{lemma}[A dyadic sum]\label{amb-shells}
Let $S$ be a set, $d_v:S\to[0,\infty)$, and
$w:S\to[0,\infty)$, with $f\ge1$ indices $v$. Suppose
\begin{align*}
 \#\{s:d_v(s)\le r_v\text{ for all }v\}
                 &\le A\exp(\kappa\sum_vr_v),\qquad r_v\ge0,\\
 w(s)&\le\prod_v\exp\left[-\frac h8
                              \min(d_v(s)^2,d_v(s)/c)\right],
\end{align*}
where the counted sets are finite and $A,\kappa\ge0$.
Assume
\begin{equation}\label{amb-shell-hypotheses}
 0<\epsilon\le1/c,\quad c>0,\quad
 \kappa\epsilon\le\log2,\quad
 32\kappa\le h\epsilon,\quad
                     T:=h\epsilon^2/16\ge1.
\end{equation}
Then the sum over $s$ for which some $d_v(s)>\epsilon$ is
absolutely convergent and at most
\begin{equation}\label{amb-shell-bound}
                          2Af3^f e^{-T}.
\end{equation}
\end{lemma}
\begin{proof}
Partition each coordinate into $[0,\epsilon]$ and
$(2^{j-1}\epsilon,2^j\epsilon]$, $j\ge1$.
For $d\ge\epsilon$ the inequality $\epsilon\le1/c$ gives
$\min(d^2,d/c)\ge\epsilon d$. Thus the weight in the $j$th
interval, $j\ge1$, is at most $e^{-T2^j}$; in interval zero it
is at most one. The number of points in a product of intervals
is bounded using the encompassing box. Set
\[
 W_0=e^{\kappa\epsilon},\qquad
 W_j=\exp\bigl((\kappa\epsilon-T)2^j\bigr)\quad(j\ge1).
\]
Their products, multiplied by $A$, bound the contributions of
the corresponding boxes. We have $W_0\le2$ and
$\kappa\epsilon\le T/2$. Hence
\[
       \sum_{j\ge1}W_j
           \le\sum_{i\ge0}e^{-T2^i}\le2e^{-T}\le1.
\]
For the middle inequality use $2^i\ge1+i$ and $T\ge1$, so
$e^{-T2^i}\le e^{-T}2^{-i}$. Write the last sum as $\beta$
and $W_0$ as $\alpha$. Summing over all multi-indices except
zero gives at most
\[
 A\bigl((\alpha+\beta)^f-\alpha^f\bigr)
       \le Af\beta(\alpha+\beta)^{f-1}
       \le2Af3^f e^{-T}.
\]
Every box is finite, so the partition also proves that the
index set is countable. The finite bounds and the convergent
nonnegative majorant justify the sum and its convergence.
\end{proof}

\begin{lemma}\label{amb-positive-type}
For the test function in \eqref{amb-idempotent},
\begin{equation}\label{amb-test-bound}
                         e(1)\ge0,\qquad |e(y)|\le e(1).
\end{equation}
If $e(1)=0$, then $e$ is identically zero.
\end{lemma}
\begin{proof}
The one-point positive-type condition makes $e(1)$ real and
nonnegative. The two-point condition says that
$\begin{pmatrix}e(1)&e(y)\\\overline{e(y)}&e(1)\end{pmatrix}$
is positive semidefinite. Its determinant is nonnegative, which
is \eqref{amb-test-bound}. This also proves the last assertion.
\end{proof}

With inner products conjugate linear in the first argument, set
\begin{equation}\label{amb-kernel-function}
 \Phi(g)=\Bigl(\prod_vd_v\Bigr)
          \overline{\prod_v\langle w_v,\sigma_v(g_v)w_v\rangle}
                 e(g_f),\qquad
 K(x,y)=\sum_{\gamma\in\Gamma}\Phi(x^{-1}\gamma y).
\end{equation}
In particular $\Phi(1)=m_0$. The representation-theoretic
estimate used for the next theorem is the precise assertion
\begin{equation}\label{amb-decay-input}
 |\langle w_v,\sigma_v(a)w_v\rangle|
       \le\exp\left[-\frac h8
                 \min\bigl(t(a)^2,t(a)/\sqrt{2q}\bigr)\right].
\end{equation}
It is needed for all $a\in G(F_v)$ and the chosen lowest-type
vectors. We now derive it, with a stronger constant, from a
classical integral formula for the lowest compact type.

\subsection{Deriving decay from the lowest compact-type formula}

We first give the elementary estimate for the scalar integrand.
We then identify the exact integral formula in a classical
representation-theoretic theorem, including its character,
measure, and projection conventions.

For a finite real random variable $Y$ supported in $[-M,M]$,
write
\[
 \varphi(u)=\log\mathbb E(e^{uY}),\qquad
 \operatorname{Var}(Y)=\mathbb E(Y^2)-\mathbb E(Y)^2.
\]
The expectation is a finite sum with nonnegative weights of sum one.

\begin{lemma}\label{amb-variance}
Suppose $\mathbb E(Y)=0$ and $0<\theta\le1$ satisfies
$2M\theta\le1$. Then
\[
                     \log\mathbb E(e^Y)
                  \ge\frac{\theta}{\mathrm e}\operatorname{Var}(Y),
\]
where $\mathrm e=\exp(1)$.
\end{lemma}
\begin{proof}
Differentiation of a finite sum shows that $\varphi''(u)$ is
the variance for the probability measure whose weights are
multiplied by $e^{uY}/\mathbb E(e^{uY})$.
For $u\ge0$ each multiplier is at least $e^{-2Mu}$.
Variance is $\min_{c\in\mathbb R}\mathbb E((Y-c)^2)$.
Taking this minimum after the pointwise comparison gives
\[
                \varphi''(u)\ge e^{-2Mu}\operatorname{Var}(Y).
\]
In particular $\varphi$ is convex. Since
$\varphi(0)=\varphi'(0)=0$ and $2M\le1/\theta$,
\[
 \varphi(\theta)
   =\int_0^\theta(\theta-u)\varphi''(u)\,du
   \ge\operatorname{Var}(Y)
             \int_0^\theta(\theta-u)e^{-u/\theta}\,du
   =\frac{\theta^2}{\mathrm e}\operatorname{Var}(Y).
\]
Convexity and $\varphi(0)=0$ give
$\varphi(1)\ge\varphi(\theta)/\theta$, proving the claim.
\end{proof}

Let $J$ be a diagonal signature matrix with $q\ge1$ negative
entries, and let
\[
             X=X^*,\qquad JX+XJ=0,\qquad t=\|X\|_{\rm HS}.
\]
The signs may occur in any order. For $1\le j<n$ define
\[
 \Delta_j(X)=\det\bigl((e^X)_{a b}\bigr)_{1\le a,b\le j}>0.
\]
For a strictly decreasing real tuple $\lambda$, put
\begin{equation}\label{amb-phase}
 \mathcal P_\lambda(X)=
       \sum_{j=1}^{n-1}(\lambda_j-\lambda_{j+1}+1)
                                      \log\Delta_j(X).
\end{equation}
Suppose that $\lambda_a-\lambda_b\ge h\ge0$ whenever
$a<b$ and the corresponding signs differ.

\begin{proposition}\label{amb-principal-minor-bound}
With these conventions,
\begin{equation}\label{amb-phase-bound}
 \mathcal P_\lambda(X)\ge
       \frac{h}{2\mathrm e}
                \min\left(t^2,\frac{t}{\sqrt{2q}}\right).
\end{equation}
\end{proposition}
\begin{proof}
The claim is immediate if $X=0$, so suppose $t>0$.
The relation $JXJ=-X$ makes the spectrum symmetric about zero.
After grouping the signs, $X$ has blocks
$\left(\begin{smallmatrix}0&B\\B^*&0\end{smallmatrix}\right)$.
It has at most $q$ positive eigenvalues, and the sum of their
squares is $t^2/2$. If $M$ is their sum, Cauchy--Schwarz gives
\[
                          M\le\sqrt{q/2}\,t.
\]

On $\bigwedge^j\mathbb C^n$, let $D_j(X)$ act by applying
$X$ to each factor and adding. Its eigenvalues are sums of
$j$ eigenvalues of $X$, hence lie in $[-M,M]$.
For the unit vector $v_j=e_1\wedge\cdots\wedge e_j$,
the identity $\bigwedge^j(e^X)=e^{D_j(X)}$ gives
\[
 \Delta_j(X)=\langle v_j,e^{D_j(X)}v_j\rangle.
\]
Expand $v_j$ in an orthonormal eigenbasis of $D_j(X)$.
The squared absolute values of its coefficients define a
finite probability measure, and its exponential moment is
the last expression. Its mean is
$\sum_{a\le j}X_{aa}=0$. Its variance is
\[
 Q_j(X)=\|D_j(X)v_j\|^2
                  =\sum_{a\le j<b}|X_{ba}|^2.
\]
Indeed the nonzero replacements of an occupied basis vector
by an unoccupied one are distinct orthogonal exterior basis
vectors; the diagonal contribution is zero.

Choose
$\theta=\min(1,1/(\sqrt{2q}\,t))$.
Then $2M\theta\le1$, so Lemma~\ref{amb-variance} gives
$\log\Delta_j(X)\ge\theta Q_j(X)/\mathrm e$.
The coefficients in \eqref{amb-phase} are positive. Summing
and interchanging two finite sums,
\[
 \begin{split}
 \sum_{j=1}^{n-1}(\lambda_j-\lambda_{j+1}+1)Q_j(X)
  &=\sum_{a<b}(\lambda_a-\lambda_b+b-a)|X_{ba}|^2\\
  &\ge h\sum_{a<b}|X_{ba}|^2
   =\frac h2\,t^2.
 \end{split}
\]
Here entries with equal signs vanish. Multiplication by
$\theta/\mathrm e$ proves \eqref{amb-phase-bound}.
\end{proof}

To state the representation-theoretic identity,
let $W$ be the lowest $K$-type in $\sigma_\lambda$,
$j:W\hookrightarrow H$ its isometric inclusion, $\tau$ its
compact representation, and $w_0$ a unit highest-weight vector.
Let $P_0$ be orthogonal projection onto $\mathbb Cw_0$ and
$d_W=\dim W$.

\begin{theorem}[The lowest compact-type integral]\label{amb-compression}
Suppose
$\lambda_1>\cdots>\lambda_n$ and
$\lambda_i-(n-1)/2\in\mathbb Z$.
Then
\begin{equation}\label{amb-compression-input}
 j^*\sigma_\lambda(e^X)j
   =d_W\int_K e^{-\mathcal P_\lambda(k^{-1}Xk)}
                 \tau(k)P_0\tau(k)^{-1}\,dk ,
 \qquad X=X^*,\ JX+XJ=0,
\end{equation}
with probability Haar measure on $K$.
\end{theorem}
\begin{proof}
We specify the two classical representation-theoretic inputs
and then prove the change of coordinates in the formula.
Let $\mathfrak b$ be the upper triangular Borel subalgebra of
$\mathfrak{gl}_n(\mathbb C)$, with strictly upper triangular
radical $\mathfrak n$, and let $T$ be the diagonal torus
of $G$. The Cartan involution is conjugation by $J$, so
$\mathfrak b$ is stable under it and its real Levi subgroup
is the compact torus $T$. Put
\[
 \rho_i=(n+1)/2-i,\qquad \nu=\lambda-\rho,\qquad
 \mu_0=\nu+2\rho_{\rm n}
                  =\lambda+\rho_{\rm n}-\rho_{\rm c}.
\]
The differences of consecutive coordinates of $\lambda$
are positive integers. Thus $\nu$ is an integral dominant
weight and defines a unitary character of $T$.
The classical realization of discrete series by
cohomological induction identifies the representation
induced from this character, denoted $A_{\mathfrak b}(\nu)$,
with $\sigma_\lambda$. Its lowest compact type has highest
weight $\mu_0$ and multiplicity one.
For the hypotheses in this identification, the admissibility
conditions of \cite[(5.1), Theorem 5.3]{vogan-zuckerman}
are exactly that $\nu$ be a unitary torus character and
$\langle \nu,e_a-e_b\rangle\ge0$ for $a<b$.
The identifying criterion
\cite[Proposition 6.1]{vogan-zuckerman}
applies to $\sigma_\lambda$: its lowest type is $\mu_0$,
and its Casimir eigenvalue is
$\|\lambda\|^2-\|\rho\|^2$. The additional condition on a
noncompact part of the Cartan subalgebra is empty because
$T$ is compact. These semisimple statements apply through
$SU(p,q)\times U(1)\longrightarrow U(p,q)$; the extra
central factor acts by the prescribed character and all
root shifts have coordinate sum zero.

We use Sun's integral formula for the lowest compact type
of $A_{\mathfrak b}(\nu)$
\cite[Theorem 6.4 and \S6.4, p.~43]{sun-eisenstein}.
In its precise specialization here, for $x\in\exp(\mathfrak p_0)$,
\begin{equation}\label{amb-sun}
 j^*\sigma_\lambda(x^2)j
 =d_W\int_K(\delta_{\mathfrak n}\nu)(\widetilde H(lx))
                       \tau(l^{-1})P_0\tau(l)\,dl.
\end{equation}
Here $\mathfrak p_0=\{X=X^*:JXJ=-X\}$ and
$\delta_{\mathfrak n}$ is the determinant of the diagonal
adjoint action on $\mathfrak n$. It has weight $2\rho$.
To define $\widetilde H$, factor $g=n a u$ in the complex
group, where $n$ is upper triangular unipotent, $a$ is
positive real diagonal, and $u\in U(n)$. Then
$\widetilde H(g)=a^2$.
Indeed Sun's definition is $y\overline{\theta(y^{-1})}$,
where the bar is conjugation with respect to the real
group $U(p,q)$; its composition with the Cartan involution
is the compact conjugation $g\mapsto(g^*)^{-1}$.
It therefore sends $a^{-1}$ to $a$.
The Haar measure in \eqref{amb-sun} is a probability
measure by the conventions of \cite[\S1.2]{sun-eisenstein}.
Sun's matrix coefficient is the inclusion followed by
the representation and projection onto the compact
isotypic component. Multiplicity one and unitarity make
it exactly $j^*\sigma_\lambda(g)j$, with value $1_W$ at
the identity. Thus no scalar factor has been suppressed.

Here is an elementary verification of the matrix conversion.
Set $x=e^{X/2}$ and $B=l e^X l^{-1}$, a positive definite
Hermitian matrix of determinant one. There is a unique
factorization
\[
                 B=n a^2 n^*,\qquad
       n\text{ upper triangular unipotent},\quad
       a=\operatorname{diag}(a_1,\ldots,a_n)>0.
\]
It follows inductively from the Schur complement of the
last diagonal entry: if
$B=\left(\begin{smallmatrix}B_0&b\\b^*&c\end{smallmatrix}\right)$,
then $c>0$ and $B_0-bb^*/c>0$; factor the latter matrix,
take $a_n^2=c$, and take the last column of $n$ above
the diagonal to be $b/c$. This also proves uniqueness.
The matrix $a^{-1}n^{-1}lx$ is unitary, so this is the
factorization used to define $\widetilde H(lx)=a^2$.

Write $z_i=a_i^2$ and
$T_i=\det B_{\{i+1,\ldots,n\},\{i+1,\ldots,n\}}$.
Triangularity gives
$T_i=\prod_{j>i}z_j$, and $\prod_j z_j=\det B=1$.
The character $\delta_{\mathfrak n}\nu$ has weight
$\kappa=\nu+2\rho=\lambda+\rho$. Summation by parts gives
\[
 \log(\delta_{\mathfrak n}\nu)(a^2)
    =\sum_j\kappa_j\log z_j
    =-\sum_{i=1}^{n-1}
             (\lambda_i-\lambda_{i+1}+1)\log T_i.
\]
The complementary-principal-minor identity, with $\det B=1$,
is
\[
       T_i=\det (B^{-1})_{\{1,\ldots,i\},\{1,\ldots,i\}}.
\]
For clarity, partition $B$ after its first $i$ rows
and columns. If its lower right block is $D$, its
Schur complement is $S$, and the upper left block of
$B^{-1}$ is $S^{-1}$. Consequently
$\det(B^{-1})_{\le i,\le i}=(\det S)^{-1}
 =\det D/\det B$, which proves the identity.
Since $B^{-1}=l e^{-X}l^{-1}$, we have established
\[
  (\delta_{\mathfrak n}\nu)(\widetilde H(lx))
                  =e^{-\mathcal P_\lambda(-lXl^{-1})}.
\]
Finally $J$ lies in the center of $K$ and $JXJ=-X$.
Replacing $l$ by $lJ$ in \eqref{amb-sun} removes the minus
sign in the phase. The projection factor is unchanged,
because $\tau(J)$ is scalar by irreducibility of $\tau$.
Replace $l$ by $k^{-1}$ to obtain
\eqref{amb-compression-input}.
\end{proof}

Thus the character shifts and dimension factor in the
source's integral formula agree with a classical theorem.
The proof uses the cited results on cohomological induction;
it does not reprove those representation-theoretic theorems.

\begin{corollary}\label{amb-decay-from-compression}
Under the hypotheses of Theorem~\ref{amb-compression},
suppose $h\ge0$ is a lower bound for all positive
opposite-sign root differences. Then for any unit
vectors $u,v$ in the lowest compact type and every $a\in G$,
\[
 |\langle u,\sigma_\lambda(a)v\rangle|
 \le\exp\left[-\frac{h}{2\mathrm e}
           \min\left(t(a)^2,\frac{t(a)}{\sqrt{2q}}\right)\right].
\]
Consequently \eqref{amb-decay-input} holds.
\end{corollary}
\begin{proof}
Compact Schur orthogonality gives
$d_W\int_K\tau(k)P_0\tau(k)^{-1}\,dk=1_W$:
the integral commutes with $\tau(K)$, and its trace is one.
The norm and the off-diagonal signature conditions on $X$
are unchanged by conjugation by $K$. Proposition
\ref{amb-principal-minor-bound} bounds every positive scalar
in \eqref{amb-compression-input} by the asserted exponential.
All the conjugates of $P_0$ are positive semidefinite.
The integral therefore lies between zero and that scalar
times $1_W$ in the order on self-adjoint operators; its norm
has the same bound.

For completeness, write $a=k e^X$ using polar decomposition.
The unitary identity gives
$J(a^*a)J=(a^*a)^{-1}$, so
$X=\frac12\log(a^*a)$ satisfies $JXJ=-X$.
The unitary factor $k$ lies in $K$, and
$\|X\|_{\rm HS}=t(a)$ by the definition using singular values.
Since $j(W)$ is $K$-stable, the compression at $a$ is
$\tau(k)$ times the compression at $e^X$ and has the same
norm. This proves the coefficient estimate.
Finally $\mathrm e<4$, for example by bounding the exponential
series by $1+1+\sum_{m\ge2}2^{-(m-1)}=3$, with strict
inequality. Thus $1/(2\mathrm e)>1/8$.
\end{proof}

The restriction to the lowest compact type is essential.
For unrestricted unit vectors in a unitary representation,
fix $a$ and a unit vector $v$, and take $u=\sigma_\lambda(a)v$.
Then $|\langle u,\sigma_\lambda(a)v\rangle|=1$.
If $h>0$ and $t(a)>0$, the displayed exponential is less
than one. Thus the corresponding assertion for all
unit vectors would be false; neither the source nor the
corollary makes that assertion.
Interchanging $J$ and $-J$ changes neither the group,
the compact subgroup, nor the phase. Hence $q$ can be
replaced by $\min(p,q)$ in the bound. In the range
$n\ge3$, $q^2\le n$ used below, one already has $q\le p$.

\subsection{The distant rational terms}

\begin{theorem}[The distant rational terms]\label{amb-far}
Assume $n\ge3$, $1\le q$, $q^2\le n$, $f\ge1$,
$r_E\ge1$, $b\ge2$, $C_g\ge0$, $1\le g\le b^{C_g}$,
$\log f\le\sqrt n$, and the lower bound for $h$ in
\eqref{amb-parameters}. Assume \eqref{amb-decay-input} and put
\begin{equation}\label{amb-truncation}
                X=nfr_Egb,\qquad \tau_0=X^{-11q}.
\end{equation}
For every $x$ with $\lambda_x\ge\tau_0$, the rational sum
\begin{equation}\label{amb-far-sum}
 \sum_{\substack{\gamma\in\Gamma\\
       t((x^{-1}\gamma x)_v)>\epsilon\text{ for some }v}}
                     |\Phi(x^{-1}\gamma x)|
                              \le m_0e^{-n^2f}
\end{equation}
is absolutely convergent, where $\epsilon$ is
\eqref{amb-small-radius}. Consequently its integral over any
measurable subset of $\{x:\lambda_x\ge\tau_0\}$ in a
finite-volume fundamental domain is at most $e^{-n^2f}D_*$.
Terms with $\Phi(x^{-1}\gamma x)=0$ may freely be included.
\end{theorem}
\begin{proof}
If $m_0=0$, Lemma~\ref{amb-positive-type} gives $e=0$, and all
assertions hold. Otherwise divide the summands by $m_0$.
Their nonzero finite components lie in $K_f$.
Lemma~\ref{amb-positive-type} and \eqref{amb-decay-input} give
the weight hypothesis of Lemma~\ref{amb-shells}, while
Corollary~\ref{amb-rational-ball} gives its counting hypothesis
with $A=A_x$, $\kappa=2n^2$ and $c=\sqrt{2q}$.

Here are all the numerical estimates needed to apply that lemma
and obtain its asserted strength. Put $L=\log(br_E)$ and
$s=\sqrt n$. We use $L\ge\log2\ge2/3$; the last inequality
follows by integrating $2/(1-u^2)\ge2$ from $0$ to $1/3$.
The lower bound for $h$ gives exactly
\begin{equation}\label{amb-T-lower}
 T=\frac{h\epsilon^2}{16}
                      \ge50(C_g+4)n^3fL.
\end{equation}
It also gives $h\epsilon\ge64n^2$ and $T\ge1$.
For example
$h\epsilon\ge8\cdot10^8(C_g+4)n^{11}f^3L$, which is
larger than $64n^2$ under the stated assumptions.
Moreover $\epsilon\le1/(2q)\le1/\sqrt{2q}$ and
$2n^2\epsilon\le2\cdot10^{-6}\le\log2$.
These prove \eqref{amb-shell-hypotheses}.

Since $X\ge2n\ge6$, we have $\tau_0<1/2$; therefore
$\lambda_x\ge\tau_0$ implies $\tau_x\ge\tau_0$.
The elementary bound $\log n\le\sqrt n$ follows from
$\log(s/2)\le s/2-1$ and $\log2\le1$. Thus
\[
 \log X\le2s+(C_g+1)L,\qquad q\le s,
 \qquad \log(2f3^f)\le4f.
\]
Using $\log3\le2$, $s\le n$ and $n\ge3$ now gives
\begin{align*}
 \log A_x+\log(2f3^f)+n^2f
 &\le2n^2f\bigl(\log3+\tfrac12\log n+22q\log X\bigr)
                                                        +4f+n^2f\\
 &\le n^2f\bigl(5+s+88n+44s(C_g+1)L+4/n^2\bigr)\\
 &\le n^3f\bigl(92+44(C_g+1)L\bigr)\\
 &\le n^3f(44C_g+182)L\\
 &\le50(C_g+4)n^3fL\le T.
\end{align*}
In the third line use $5+4/n^2\le3n$; in the fourth use
$92\le138L$. Lemma~\ref{amb-shells} therefore bounds the
normalized sum by $A_x2f3^fe^{-T}\le e^{-n^2f}$.
Multiplication by $m_0$ proves \eqref{amb-far-sum}.
The sum is measurable, since it is a countable sum of measurable
nonnegative terms. Integrating the pointwise bound and using
the volume of the full domain proves the final assertion.
\end{proof}

\subsection{The trace identity and the remaining comparison}

The following statement records the exact analytic conclusion
needed from the automorphic-kernel construction. It is not a
consequence of the lattice packing argument: the convolution
operator $R(\Phi)$ must be a finite-rank orthogonal projection
of rank \eqref{amb-D-definition}, and its continuous kernel must
have diagonal
\begin{equation}\label{amb-projection-diagonal}
                  K(x,x)=\sum_{j=1}^D|F_j(x)|^2
\end{equation}
for an orthonormal basis $F_1,\ldots,F_D$ of its image, with
the series in \eqref{amb-kernel-function} representing this
kernel. With these assertions supplied, finite summation and
orthonormality immediately give
\begin{equation}\label{amb-trace}
                 K(x,x)\ge0,\qquad
                 \int_{\mathcal X}K(x,x)\,dx=D.
\end{equation}
In particular, the diagonal is integrable. A finite-rank kernel
on $L^2$ must be given by such representatives, or an equivalent
proved diagonal identity; equality of measurable kernels only
almost everywhere on $\mathcal X\times\mathcal X$ would not
by itself determine the diagonal.

\begin{proposition}\label{amb-final-reduction}
Assume $n\ge e^{15}$ and $f\ge1$. Suppose the actual kernel
satisfies \eqref{amb-projection-diagonal} and its diagonal
series is absolutely convergent. Let
\[
 \mathcal M=\{x:\lambda_x\ge\tau_0\},\qquad
 \mathcal B=\mathcal X\setminus\mathcal M,
 \qquad a=nf\log n,\qquad u=n^2f.
\]
Assume the three further estimates
\begin{align}
 \int_{\mathcal M}
   \sum_{\substack{\gamma\ne1\\
          t((x^{-1}\gamma x)_v)\le\epsilon\ (\forall v)}}
                      |\Phi(x^{-1}\gamma x)|\,dx
                  &\le e^{-a/4}D_*,                    \label{amb-near-input}\\
 \operatorname{vol}(\mathcal B)
                  &\le e^{-a/4}\operatorname{vol}(\mathcal X),
                                                        \label{amb-rare-input}\\
 \int_{\mathcal B}K(x,x)\,dx&\le e^{-u}D.             \label{amb-deep-input}
\end{align}
Under the hypotheses of Theorem~\ref{amb-far}, these assertions
imply \eqref{amb-target}.
\end{proposition}
\begin{proof}
On $\mathcal M$, separate the identity term $m_0$, the nonidentity
terms with all distances at most $\epsilon$, and the remaining
terms. Absolute convergence, \eqref{amb-near-input} and
Theorem~\ref{amb-far} give
\[
 \left|\int_{\mathcal M}K(x,x)\,dx
                      -m_0\operatorname{vol}(\mathcal M)\right|
                  \le(e^{-a/4}+e^{-u})D_*.
\]
On $\mathcal B$, positivity in \eqref{amb-trace} gives
\[
 \left|\int_{\mathcal B}K(x,x)\,dx
                      -m_0\operatorname{vol}(\mathcal B)\right|
 \le e^{-u}D+m_0\operatorname{vol}(\mathcal B).
\]
Adding and using \eqref{amb-rare-input} proves the exact inequality
\begin{equation}\label{amb-error-before-absorption}
             |D-D_*|\le(2e^{-a/4}+e^{-u})D_*+e^{-u}D.
\end{equation}
Set $z=e^{-a/16}$ and $Y=|D-D_*|$. Since $\log n\le n$,
$u\ge a$, and $n\ge e^{15}$ gives $a\ge240$.
Consequently
\[
        e^{-a/4}\le z/8,\qquad e^{-u}\le z/8,\qquad z\le1;
\]
for the first two bounds it is enough that
$3a/16\ge\log8$ and $15a/16\ge\log8$.
Use $D\le D_*+Y$ in \eqref{amb-error-before-absorption}.
It follows that
\[
                    Y\le(z/2)D_*+Y/2,
\]
whence $Y\le zD_*$. This is \eqref{amb-target}, including the
case $D_*=0$.
\end{proof}

\subsection{Hypothesis checks and precise scope}

The small-distance condition does not force a rational element
to be the identity in the indefinite group. For example, in a
hyperbolic plane with Hermitian matrix
$\begin{pmatrix}0&1\\1&0\end{pmatrix}$, choose a nonzero
$a\in\mathcal O_E$ with $\bar a=-a$. Then
\[
 \gamma=\begin{pmatrix}1&Pa\\0&1\end{pmatrix},\qquad
 x_R=\begin{pmatrix}R&0\\0&R^{-1}\end{pmatrix}
\]
are isometries, $\gamma$ is a nonidentity integral unipotent
congruent to one modulo $P$, and
$x_R^{-1}\gamma x_R=\begin{pmatrix}1&Pa/R^2\\0&1\end{pmatrix}$
tends to one at every infinite place as $R\to\infty$.
One can take the same $R$ at every place and adjoin a positive
orthogonal summand of any dimension. This satisfies the actual
unimodularity and congruence conditions. It explains why
\eqref{amb-near-input} is a separate counting problem, even
after Proposition~\ref{amb-near-unipotent} is proved.

If the congruence condition is omitted, $\gamma=-1$ has distance
zero at every place and is not unipotent. If the small-distance
condition is omitted, the product
\[
 \begin{pmatrix}1&Pa\\0&1\end{pmatrix}
 \begin{pmatrix}1&0\\Pa&1\end{pmatrix}
 =\begin{pmatrix}1+P^2a^2&Pa\\Pa&1\end{pmatrix}
\]
is an integral isometry congruent to one modulo $P$, but has
trace $2+P^2a^2\ne2$, so is not unipotent. These are
counterexamples to weakened statements, not to the actual
propositions above. In a definite group a unitary unipotent is
the identity, by the spectral theorem. Thus, once the actual
kernel and trace identity are supplied, the definite case gives
the exact equality $D=D_*$.

The present section proves the rational unipotence and separation
arguments, the orbit packing, the full distant-term estimate
using the classical integral formula and its elementary
decay estimate, and the exact deduction
of \eqref{amb-target} from the displayed analytic inputs.
It retains the explicitly cited classical representation
theorems. It does not supply the automorphic
finite-rank kernel construction, the integrated nonidentity
unipotent estimate \eqref{amb-near-input}, the small-volume bound
\eqref{amb-rare-input}, or the cusp estimate
\eqref{amb-deep-input}. These are separate substantive inputs
in the actual source chain. No counterexample satisfying the
actual hypotheses was found in the components proved here.
The Galois-cohomology bounds at other coefficient primes are a
different dependency and are not consequences of this complex
multiplicity argument.


